% Ad Familiares 9.13 (ad-familiares-09-13) --- English translation
% Translated by Claude (Anthropic) from the Latin (Perseus).
% Style guide: STYLE.md.

\ciceroLetterOpener{CICERO DOLABELLAE S.}

\ciceroSection{1} C. Subernius of Cales is both a close friend of mine and a very intimate connection of our Lepta. To escape the war, he had set out for Spain with M. Varro before the war began, expecting to be in a province where, after the defeat of Afranius, none of us thought there would be any further war at all; and there he fell into the very evils he had so anxiously gone to avoid. He was overtaken by a sudden war --- the war that was stirred up by Scapula and was afterwards so consolidated by Pompeius that there was no means by which he could extricate himself from that miserable plight.

\ciceroSection{2} Much the same is the case of the heir of M. Planius, who is likewise a man of Cales and likewise a most intimate connection of our Lepta. These two men, then, I commend to you in such terms that I could not commend them with greater care, with greater concern, with greater anxiety of mind. I want it for their own sake, and in this both friendship and humanity move me powerfully; but Lepta is so distressed --- for his own fortunes seem to be on the brink of ruin in the matter --- that I cannot help being distressed myself, next to him if not equally with him. And so, although I have often had proof of how much you love me, still I should like you to be persuaded that on this matter I shall most of all take the measure of it.

\ciceroSection{3} I therefore ask you --- or, if you will allow me, beg you --- to preserve safe and sound men who are wretched and undone more by a fortune that no one can avoid than by any fault of their own, and to be willing that this gift, through your agency, should be one I confer at once on the men themselves --- friends of mine --- and on the township of Cales, with which I have a great bond, and above all on Lepta, whom I put before everyone.

\ciceroSection{4} What I am about to say is, I admit, not greatly to the point; but still no harm in saying it: the family property of one of them is very slight, and that of the other barely of equestrian rank. Therefore, since Caesar in his generosity has granted them their lives, and since there is no great property left for them to be deprived of besides, secure for them their return home, if you love me as much as you certainly do love me --- a return in which the only obstacle is the length of the journey, which they do not shrink from for this reason: that they may both live among their own people and die at home. That you should strive and contend for this --- or rather that you should accomplish it (for that you have the power I have persuaded myself) --- I beg you again and again, with all my heart.
